\chapter{讨论与结论}

\section{对案例研究的启示}

这一公开示例表明，职业转型研究完全可以在不接触任何隐私数据的前提下完成完整写作训练。对于模板仓库而言，这一点尤其重要：示例项目的任务不是替代真实论文，而是展示“结构、样式与构建链路”如何在公开数据条件下稳定运行。

\section{方法局限}

本示例采用的是公开案例研究，因此结论不应被解释为行业普遍规律。它更适合作为“结构化写作演示”：帮助使用者理解如何组织摘要、章节、图表、长表格与参考文献，而不是给出关于娱乐产业的最终经验判断。

\section{结论}

本研究形成三点结论。第一，跨媒介职业转型具有明显的阶段性，可用出道、扩张、危机与重启四个阶段进行概括。第二，平台自主性是危机后形象恢复的重要条件。第三，公开案例足以支撑一个内容完整、版式丰富且可复现构建的毕业论文示例项目。

\section{模板使用建议}

对后续使用者而言，推荐遵循三条原则：
\begin{enumerate}
  \item 把元信息集中维护在 `extraTex/meta.tex`；
  \item 把正文按章节拆分到 `extraTex/body/`；
  \item 在保留真实版式的前提下，优先使用公开可分享的示例文字与图表。
\end{enumerate}
